\documentclass[11pt,a4paper]{article}
\usepackage[margin=2.5cm]{geometry}
\usepackage{booktabs}
\usepackage{hyperref}
\usepackage{parskip}

\title{Hallucination Case Study\\Assignment 3 --- Smart Water Lab}
\author{Mahmudul Hasan (4125999049)}
\date{2026}

\begin{document}
\maketitle

\section{AI claim}
During test design (Week 8 Session A, Cursor Agent): ``For CN=95 and rainfall $P=5$~mm, runoff is always zero because high CN means no runoff for small storms.''

\section{Why this is wrong}
SCS-CN initial abstraction: $I_a = 0.2 \times (25400/\mathrm{CN} - 254)$.\\
For CN=95: $I_a \approx 5.37$~mm, so $P=5$~mm gives $Q=0$ \emph{by arithmetic}, not by a general rule. High CN \emph{lowers} $I_a$; small storms can still produce runoff when $P > I_a$.

\section{Tests created}
\texttt{tests/test\_runoff.py}: \texttt{test\_runoff\_below\_initial\_abstraction}, \texttt{test\_runoff\_not\_exceed\_rainfall\_high\_cn}.\\
\texttt{src/validation.py}: \texttt{validate\_runoff\_mm} enforces $Q \le P$.

\section{Failure discovered}
Draft documentation over-generalized one (CN, $P$) pair. pytest + Ia table showed the verbal rule was misleading for test design.

\section{Fix applied}
\begin{itemize}
  \item Zero-runoff test: $P=2$~mm, CN=80 ($I_a \approx 20.3$~mm).
  \item Physical sweep: CN = 60, 80, 95 at $P=100$~mm with $Q \le P$ assert.
  \item Documented in \texttt{prompt\_log.md} and Week 8A lab report.
\end{itemize}

\section{Evidence checklist}
\begin{tabular}{@{}ll@{}}
\toprule
Step & Location \\
\midrule
AI prompt & \texttt{prompt\_log.md} Week 8A \\
Tests & \texttt{tests/test\_runoff.py} \\
Failure documented & This file + Week 8A PDF \\
Fix & Corrected tests + validation layer \\
Run & \texttt{pytest -q} $\rightarrow$ 29 passed \\
\bottomrule
\end{tabular}

\end{document}
